% =====================================================================
\subsection{Verificación del sistema integrado}
\label{sec:calidad}
\label{sec:calidad-resumen}
\label{sec:control-calidad}
\label{sec:defectos}

\textbf{Pregunta.} Con cada etapa evaluada por separado, ¿el sistema ya integrado funciona, resiste
carga y protege los datos de quien lo usa?

\textbf{Método.} Es la parte de pruebas del OE5 (fase F5). El equipo de calidad ejecutó un plan de diez
fases, ordenado por las características del modelo de calidad ISO/IEC 25010 \citep{iso_25010_2011}: 0 línea
base, 1 escenarios de requisitos no funcionales, 2 unitarias y cobertura, 3 integración, 4 funcionales,
5 adversariales de IA, 6 patrones de resiliencia, 7 caos, 8 carga y 9 telemetría. Las fases 0 a 6 se corrieron
en octubre y la campaña de caos, carga y telemetría llegó hasta el 5 de noviembre de 2026.

\textbf{Métrica.} Pruebas aprobadas, cobertura de sentencias, tiempo de recuperación, errores y percentil 95 de
latencia bajo carga, y hallazgos por severidad. \textbf{Datos.} Los de cada nivel, resumidos en la
Tabla~\ref{tab:niveles-prueba}, más dos auditorías de código: la de autorización (152 puntos de entrada) y la
de datos personales conforme a la Ley 1581 de 2012 \citep{congreso_ley1581_2012}.

\begin{table}[htbp]
\centering
\small
\caption{Niveles de prueba del sistema integrado.}
\label{tab:niveles-prueba}
\begin{tabularx}{\textwidth}{@{}p{2.3cm}XXX@{}}
\toprule
Nivel & Qué verifica & Datos & Resultado \\
\midrule
Unitarias y cobertura & Bibliotecas y servicios, con los defectos conocidos marcados como fallo esperado & 2.542 pruebas Python, 86 de la web y 67 de contrato & Todas pasan; cobertura de 54\,\% (catálogo) a 98\,\% (puerta de enlace) \\
Integración & Los servicios contra el sistema desplegado completo & Bases de prueba aisladas y fronteras entre servicios & Pasan; la prueba de humo pasó \\
Funcionales & Percepción y emparejamiento con conjuntos de oro & G01-G10, lista base de 18 renglones; el oro fija la respuesta del modelo y mide al sistema & 0 productos fuera del catálogo \\
Adversariales de IA & OWASP LLM 2025 & 143 pruebas & 120 verdes; inyección en texto e imagen y producto inexistente cumplen \\
Resiliencia y caos & Caídas inyectadas con Toxiproxy & Campaña del 5 de noviembre: 25 ensayos & 14 cumplen \\
Carga & Rampa de 5 a 100 usuarios (L-02) & Apple M4 y Windows con WSL2 & 0 errores en el M4; saturación en WSL2 \\
Telemetría & Expedientes por identificador de correlación & Trazas de 2 expedientes (32 tramos) & No cumple \\
Seguridad & Pago sin datos de tarjeta; autorización & Barrido de 248 columnas & 0 números de tarjeta \\
\bottomrule
\end{tabularx}
\end{table}

\textbf{Resultado: pruebas automáticas.} De las 2.542 pruebas Python ninguna falló y una quedó saltada
(la prueba en vivo del asistente, que necesita la llave del modelo). La integración continua no exige un
mínimo de cobertura, así que esta describe el estado y no una meta.

\textbf{Resultado: adversariales de IA.} De 143 pruebas, 23 documentan un defecto. Dejaron hallazgos el
esquema roto (cuatro formas devolvían error 500 en lugar de 422), la fuga del prompt, los datos personales, la
agencia excesiva, los embeddings, el consumo sin tope y los identificadores de correlación hostiles.

\textbf{Resultado: requisitos no funcionales.} La matriz tiene 52 requisitos en 99 filas: 54 ejecutadas,
25 por medir, 9 omitidas y 8 no medibles. 

\textbf{Resultado: degradación.} Con el catálogo colgado, el disyuntor abre a los 26,3~s, dentro de los 30~s
de ESC-R2. La campaña del 5 de noviembre dio 26,4 y 26,3~s.

\textbf{Resultado: caídas.} En la campaña de octubre fallaron tres escenarios; su corrección se verificó en una rama aparte, que la versión evaluada no integra. En CH-04, con el catálogo caído,
la subida de la lista devolvía error 500 y no se guardaba. En CH-12 y CH-15, con PostgreSQL o RabbitMQ caídos,
los procesos de fondo no volvían, y un pago avisaba dos veces. Tras la corrección (verificación del 20 de
octubre), la subida responde 202 en 33~ms y se recupera en 2,7~s, CH-12 en 7,1~s, CH-15 en 6,9~s, y el pago manda un solo aviso.

\textbf{Resultado: carga.} El requisito RNF05 (decisión de diseño DP-02) pide 50 usuarios con menos de
5\,\% de error. En un Apple M4 con los servicios nativos (5 de noviembre), la rampa completó 25.706 peticiones
sin un error y sin saturación hasta 100 usuarios; el percentil 95 del catálogo (la demora que no supera el
95\,\% de las peticiones) fue de 49 a 420~ms hasta 50 usuarios y de 2,5~s con 100. En Windows con WSL2 y
4 CPU, el catálogo se saturó desde 30 usuarios, con un núcleo al 100\,\%. Por eso RNF05 se cumple en el M4,
con el doble de usuarios pedido. El umbral adicional de 300~ms para el percentil 95 se supera desde 25 usuarios
porque el catálogo corre en un solo proceso (Sección~\ref{sec:trabajo-futuro}).

\textbf{Resultado: seguridad.} La auditoría de dependencias de la web quedó en cero vulnerabilidades. Si
alguien pide el kit de otra persona, el sistema responde con error 401 (no autenticado) cuando no hay
sesión y con error 404 (no existe) cuando la sesión es de otra cuenta. La auditoría de autorización halló
22 de 152 puntos de entrada que no comparaban por completo a qué cuenta pertenece cada dato, y de seis
hallazgos de esa revisión se cerraron cuatro. El defecto de mayor alcance, las políticas de acceso a nivel
de fila de la base de datos que concedían la tabla completa, se evita porque el sistema controla ese acceso
en la aplicación (Sección~\ref{sec:autorizacion}).

\textbf{Resultado: hallazgos de calidad.} El registro del control de calidad cuenta 85 hallazgos vigentes del sistema evaluado
al 5 de noviembre (16 de severidad alta, 58 media y 11 baja) y 10 cerrados. Los defectos de código tienen además una prueba que se espera
que falle (\texttt{xfail} estricto): si el defecto se corrige y la prueba pasa, la prueba falla y obliga
a cerrar el hallazgo. Entre los de severidad alta están las rutas de administración sin sesión, el
administrador sin segundo factor de autenticación y el nombre del menor, que llega sin depurar al
proveedor del modelo.

El recorrido completo se evalúa así. Con pago simulado, dos corridas
automáticas dieron resultados distintos: el pedido llegó a «pagado» el 1 de noviembre y quedó «pendiente
de pago» el 2, por un fallo de la consulta de estado que no se repitió. En cambio, en las seis sesiones
con acudientes el pedido terminó pagado. Con eso se da por funcional el recorrido hasta el pago simulado.

La Tabla~\ref{tab:iso25010} ordena la evidencia por característica de la ISO/IEC 25010 («Parcial»: falta
alguna prueba o hay hallazgos abiertos). T4, T6 y C3 se definen en la Sección~\ref{sec:oe5}.

\begin{table}[htbp]
\centering
\scriptsize
\caption{Evidencia de calidad por característica de la ISO/IEC 25010.}
\label{tab:iso25010}
\begin{tabularx}{\textwidth}{@{}lXl@{}}
\toprule
Característica & Evidencia & Estado \\
\midrule
Adecuación funcional & 2.542 pruebas Python y las de las dos aplicaciones web; recorrido de la cuenta al pago con acudientes; personaje del segundo hijo sin reflejar en el kit Fan & Parcial \\
Eficiencia de desempeño & ESC-D4 con p95 de 0,31~s frente a 8~s; L-02 hasta 100 usuarios sin errores en el M4, con p95 del catálogo de 49 a 420~ms hasta 50 usuarios, y saturación desde 30 en la máquina virtual & Parcial \\
Compatibilidad & Pruebas de fronteras entre servicios (75) y del contrato entre la web y la API (67) & Probada \\
Usabilidad & Seis sesiones con acudientes: T6 4 de 5 y T4 1 de 5 sin ayuda; SUS medio de 90 [82--98]; C3 solo en T6 & Parcial \\
Fiabilidad & Campaña de caos del 5 de noviembre (14 de 25 ensayos cumplidos); hallazgos de idempotencia y avisos abiertos & Parcial \\
Seguridad & Lectura de kits ajenos cerrada, 0 vulnerabilidades web, 0 de 81 ataques al asistente; un token en el historial; hallazgos de severidad alta abiertos & Parcial \\
Mantenibilidad & Cobertura de 54 a 98\,\% por servicio, sin umbral en la integración continua & Parcial \\
Portabilidad & Infraestructura con Docker Compose y servicios locales con el sistema desplegado completo; sin pruebas en dispositivos & Parcial \\
\bottomrule
\end{tabularx}
\end{table}

\textbf{Lectura.} El sistema integrado cumple la capacidad pedida y el recorrido hasta el pago simulado
funciona; sin embargo, solo una de las ocho características de la ISO/IEC 25010 quedó probada por
completo. Los hallazgos de seguridad y privacidad deben cerrarse antes de operar con datos y pagos
reales, y por eso encabezan el trabajo futuro.
